\chapter{Datos y construcción del grafo vial}\label{cap:grafo}

Este capítulo explica de dónde salen las carreteras, cómo se convierten en un grafo utilizable por los algoritmos y qué decisiones de diseño se tomaron en el
camino, incluidas las que no funcionaron.

\section{Fuente de los datos}

Las carreteras proceden de \emph{OpenStreetMap} (OSM)~\cite{haklay2008}, una base de datos geográfica colaborativa publicada bajo la licencia ODbL (\emph{Open Database License}), que obliga a citar la
fuente como «\textcopyright\ OpenStreetMap contributors». No es una fuente oficial: la calidad de los datos depende de la comunidad de mapeadores, lo que implica
que algunas etiquetas (por ejemplo, la velocidad máxima de un tramo) están incompletas. Aceptamos esta limitación porque OSM es la única fuente abierta que describe toda la red
con la topología necesaria (tramos, sentidos de circulación y jerarquía de las vías).

Usamos el extracto de España de \emph{Geofabrik} (\texttt{spain-latest.osm.pbf}, 1,4~GB, última modificación el 19 de octubre de 2026). Para las ciudades empleamos datos oficiales del INE
(capítulo~\ref{cap:versiones}).

\section{Un primer enfoque descartado: la API Overpass}

El primer enfoque fue descargar la red directamente con OSMnx~\cite{boeing2017} a través de la API pública \emph{Overpass}, dividiendo España en 324 consultas. En unas once horas solo se completaron
19 consultas (el 6~\%), por errores \texttt{ConnectTimeout} y \texttt{502 Bad Gateway} del servidor público, que estaba saturado; además, una de sus dos direcciones IP no aceptaba conexiones y
OSMnx fija una sola IP por servidor. Estimamos que completar la descarga llevaría días y no sería reproducible. Se conserva el registro de aquel intento como evidencia
(\texttt{docs/intento\_overpass/}), pero se descartó.

\recuadro{Decisión}{Sustituir Overpass por un único archivo descargable y reanudable (el extracto de Geofabrik) y filtrarlo en local. La descarga completa tarda unos siete minutos y puede repetirse.}

\section{Filtrado y construcción del grafo}\label{sec:construccion}

El extracto de 1,4~GB contiene todos los elementos de OSM de España. Nos interesa solo la red de carreteras de coche, así que conservamos las vías cuya etiqueta
\texttt{highway} es \texttt{motorway} (autopista y autovía), \texttt{trunk} (vía rápida), \texttt{primary} (red principal) o \texttt{secondary} (red secundaria), junto con sus
enlaces (\texttt{*\_link}). El filtrado se hace con la biblioteca \texttt{pyosmium} en dos pasadas: la primera selecciona las vías que cumplen el filtro; la segunda recupera las coordenadas de los nodos que esas vías usan.

El grafo resultante es un \emph{multigrafo dirigido}: cada tramo de carretera de doble sentido se representa con dos aristas, una en cada dirección, y una autovía con
calzadas separadas son dos aristas distintas. Se admite más de una arista entre el mismo par de nodos (carreteras paralelas), que se distinguen por una clave.

\subsection{Simplificación}

Un grafo con todos los puntos del trazado tendría millones de nodos. Se \emph{simplifica}: solo se conservan como nodos las intersecciones y los extremos de carretera, y cada arista
representa un tramo completo entre dos de ellos, con su longitud total. Seguimos las mismas reglas que OSMnx: un nodo se conserva si tiene bucle sobre sí mismo, si no tiene aristas
de entrada o de salida, o si no es un simple punto intermedio (un nodo con exactamente dos vecinos y grado 2 o 4 se considera intermedio y se fusiona). La longitud de un tramo es la suma de los segmentos
originales, medidos con la distancia ortodrómica (\emph{great-circle}) entre coordenadas.

\subsection{Conectividad}

Para que exista un camino entre cualquier par de nodos, se conserva la mayor componente \emph{fuertemente conexa} del grafo, es decir, aquella donde se puede ir de cualquier nodo a cualquier otro respetando
los sentidos de circulación. Esto elimina también los fragmentos aislados (islas, tramos desconectados, extremos de frontera). El grafo base final es fuertemente conexo y cubre únicamente la
\textbf{España peninsular}: Baleares, Canarias, Ceuta y Melilla no están.

\section{Atributos del grafo}

Cada arista guarda tres valores numéricos y varios descriptivos. La velocidad se toma de la etiqueta \texttt{maxspeed} de OSM cuando existe; si falta, se imputa la media de los tramos
de su mismo tipo de vía. Con la longitud $L$ en metros y la velocidad $v$ en km/h, el tiempo de viaje en segundos es
\begin{equation}\label{eq:tiempo}
  t\;=\;\frac{L}{v/3{,}6}.
\end{equation}
Por ejemplo, un tramo de la AP-8 de $5.895{,}57$~m a $100$~km/h tiene un tiempo de $212{,}24$~s. La tabla~\ref{tab:atributos} resume los atributos.

\begin{table}[htbp]
  \centering
  \caption{Atributos de las aristas del grafo base.}\label{tab:atributos}
  \small
  \begin{tabularx}{\linewidth}{@{}l l c X@{}}
    \toprule
    Atributo & Tipo & Presencia & Significado \\
    \midrule
    \texttt{length} & decimal & 100 \% & Longitud del tramo en metros (peso para distancias). \\
    \texttt{travel\_time} & decimal & 100 \% & Tiempo de viaje en segundos, según la ecuación~\eqref{eq:tiempo} (peso para tiempos). \\
    \texttt{speed\_kph} & decimal & 100 \% & Velocidad usada (km/h): la máxima de OSM o la media de su tipo de vía. \\
    \texttt{highway} & texto & 100 \% & Tipo de vía (\texttt{motorway}, \texttt{trunk}, \texttt{primary}, \texttt{secondary} o sus enlaces). \\
    \texttt{oneway} & booleano & 100 \% & \texttt{True} si es de sentido único (227.446 aristas) o \texttt{False} si es de doble sentido (81.804). \\
    \texttt{ref} & texto & 90,3 \% de los km & Código de la carretera (\texttt{AP-7}, \texttt{A-4}, \texttt{N-340}\dots). \\
    \texttt{name} & texto & 72,0 \% de los km & Nombre en lenguaje natural («Autovía del Mediterráneo»). \\
    \bottomrule
  \end{tabularx}
\end{table}

Los nodos guardan su longitud (\texttt{x}), su latitud (\texttt{y}) y el número de calles que confluyen (\texttt{street\_count}). OSM no nombra las intersecciones, por lo que
\emph{los nodos no tienen nombre}: los nombres de ciudad se asocian a un nodo en el capítulo~\ref{cap:versiones}.

Para poder guardar el grafo en un formato estándar se eliminaron dos atributos de OSM: \texttt{osmid} (listas de identificadores, inútiles para trabajar y no admitidas por el formato) y
\texttt{maxspeed} (con listas y valores no numéricos como \texttt{ES:urban} o \texttt{50 mph}, cuya información ya está resuelta en \texttt{speed\_kph}). Los atributos
que OSM da como lista (\texttt{ref}, \texttt{name}) se unieron en un único texto.

\subsection{Formato de archivo}

El grafo se guarda en \emph{GraphML}, un formato estándar de grafos basado en XML, con \texttt{nx.write\_graphml} de NetworkX. Este formato declara el tipo real de cada atributo,
de modo que al leerlo con \texttt{nx.read\_graphml} los números llegan como números y los booleanos como booleanos, y el \emph{notebook} solo necesita NetworkX, como las prácticas de la asignatura.
El archivo ocupa 108~MB y se carga en unos 20 segundos.

\section{Descripción del grafo obtenido}

El grafo base tiene \textbf{179.594 nodos} y \textbf{309.250 aristas dirigidas} (unos 183.000~km contando cada sentido). Hay 3.349 pares de nodos unidos por más de una arista.
La tabla~\ref{tab:tipos} reparte las aristas por tipo de vía y la tabla~\ref{tab:nodos} clasifica los nodos según el número de calles que confluyen.

\begin{table}[htbp]
  \centering
  \caption{Aristas y kilómetros por tipo de vía.}\label{tab:tipos}
  \small
  \begin{tabularx}{0.92\linewidth}{@{}l X r r@{}}
    \toprule
    Tipo (\texttt{highway}) & Qué es & Aristas & Km \\
    \midrule
    \texttt{motorway} & autopista y autovía & 17.618 & 31.485 \\
    \texttt{trunk} & vía rápida o carretera desdoblada & 52.029 & 31.438 \\
    \texttt{primary} & carretera de la red principal & 75.467 & 40.454 \\
    \texttt{secondary} & carretera de la red secundaria & 100.495 & 69.150 \\
    \texttt{*\_link} & enlaces y accesos de los cuatro anteriores & 63.641 & 10.233 \\
    \midrule
    \multicolumn{2}{@{}l}{\textbf{Total}} & \textbf{309.250} & \textbf{182.760} \\
    \bottomrule
  \end{tabularx}
\end{table}

\begin{table}[htbp]
  \centering
  \caption{Tipos de nodo según el número de calles que confluyen (\texttt{street\_count}).}\label{tab:nodos}
  \small
  \begin{tabularx}{0.92\linewidth}{@{}c X r@{}}
    \toprule
    Calles & Significado típico & Nodos \\
    \midrule
    1 & extremos (fin de carretera o corte en la frontera) & 1.058 \\
    2 & empalmes entre dos vías distintas de OSM & 4.612 \\
    3 & bifurcaciones o incorporaciones (entradas y salidas de autovía) & 161.516 \\
    4 & cruces de cuatro ramas & 10.332 \\
    5 o más & nudos complejos & 2.076 \\
    \bottomrule
  \end{tabularx}
\end{table}

\section{Memoria y rendimiento: reescribir la simplificación}

El grafo con carreteras secundarias no cabía en memoria con el flujo habitual de OSMnx (\texttt{graph\_from\_*}): se estimaban unos 12~GB y el proceso fue finalizado por el sistema operativo por falta
de memoria. En lugar de eso, el Anexo~B reimplementa con \emph{arrays} de NumPy las reglas de simplificación de OSMnx, sin construir un grafo gigante de NetworkX intermedio. El pico de memoria medido es de
unos 2,2~GB y la construcción completa tarda unos cuatro minutos.

Para asegurarnos de que la reimplementación es fiel, la comparamos con el grafo \emph{ligero} (solo \texttt{motorway}, \texttt{trunk}, \texttt{primary} y sus enlaces) que ya habíamos construido con el método anterior:
los dos tienen los mismos 108.050 nodos (con los mismos identificadores), las mismas 181.006 aristas y las mismas longitudes (110.128,4~km). Solo cambian ligeramente las velocidades imputadas (diferencia media de
0,05~km/h; 136 aristas difieren más de 1~km/h), porque el método anterior elegía de forma arbitraria el tipo de vía de las aristas mixtas y el nuevo usa el de mayor jerarquía.

\section{Por qué incluir las carreteras secundarias}

El grafo base incluye carreteras secundarias; el ligero no. Con solo autopistas, vías rápidas y red principal, una red \emph{sin autovías} queda llena de agujeros: Almería quedaba a 142~km de la red, y la ruta
Madrid--Sevilla medía 761~km en lugar de unos 510. Con las secundarias el resultado es utilizable, y además las ciudades vecinas quedan mucho más cerca entre sí (en la red sin secundarias, Terrassa--Sabadell salía un 193~\% más lejos).
Las versiones más pequeñas del grafo se obtienen después filtrando aristas del grafo base, así que mantener una sola red «maestra» no supone un problema de cómputo.

\section{Limitaciones}

\begin{itemize}
  \item \textbf{Sin geometría.} Al estar simplificadas, las aristas no guardan su forma: dibujadas, son líneas rectas entre intersecciones. Las distancias no se ven afectadas, pero los mapas son aproximados: las rutas
        dibujadas en la web (capítulo~\ref{cap:web}) siguen la sucesión de intersecciones del camino mínimo, no la forma exacta de cada curva.
  \item \textbf{Tiempos de coche idealizados.} \texttt{travel\_time} no incluye tráfico, peajes ni semáforos, y muchas velocidades están imputadas.
  \item \textbf{Longitud geométrica.} \texttt{length} es la longitud del trazado en OSM, no el kilometraje oficial.
  \item \textbf{Datos colaborativos.} Pueden contener errores o carencias locales.
\end{itemize}
